\section{Command}
\label{sec:pat-command}

\problemstatement{Encapsulate a request as an object, so requests can be
parameterised, queued, logged, and undone -- decoupling the object that
invokes an operation from the one that knows how to perform it.}

\begin{patternbox}[When You Reach for It]
The triggers are \emph{``undo/redo,''} \emph{``queue or schedule
operations,''} and \emph{``parameterise a button/menu with an action.''}
When an action must be treated as a first-class value -- stored, passed,
reversed -- wrap it in a Command object.
\end{patternbox}

\subsection*{Understanding the pattern}

A Command object exposes \texttt{execute()} (and often \texttt{undo()})
and holds everything the action needs: a reference to the receiver and any
arguments. An invoker (a button, a queue, a macro) triggers commands
without knowing what they do. Because each command is an object, you can
keep a history stack for undo, replay a log, or batch commands into a
macro.

\begin{ideabox}[Turn a Method Call into an Object]
Reifying ``call this method with these arguments'' as an object is what
enables undo, queuing, logging and macros -- none of which are possible
with a bare function call. The invoker depends only on the Command
interface.
\end{ideabox}

\subsection*{C++ implementation}

\begin{lstlisting}
#include <bits/stdc++.h>
using namespace std;

struct Light {                                  // receiver
    bool on = false;
    void turnOn()  { on = true;  cout << "Light on\n"; }
    void turnOff() { on = false; cout << "Light off\n"; }
};

struct Command {
    virtual void execute() = 0;
    virtual void undo() = 0;
    virtual ~Command() = default;
};

struct LightOnCommand : Command {
    Light& light;
    LightOnCommand(Light& l) : light(l) {}
    void execute() override { light.turnOn(); }
    void undo() override { light.turnOff(); }
};

class RemoteControl {                            // invoker with undo
    vector<unique_ptr<Command>> history;
public:
    void submit(unique_ptr<Command> cmd) {
        cmd->execute();
        history.push_back(move(cmd));
    }
    void undoLast() {
        if (!history.empty()) {
            history.back()->undo();
            history.pop_back();
        }
    }
};

int main() {
    Light light;
    RemoteControl remote;
    remote.submit(make_unique<LightOnCommand>(light));
    remote.undoLast();                          // reverses it
    return 0;
}
\end{lstlisting}

\begin{complexitybox}[Trade-offs]
\textbf{Pros.} Enables undo/redo, queuing, logging, and macros; decouples
invoker from receiver; commands are composable. \textbf{Cons.} A class per
action can proliferate; undo requires each command to capture enough state
to reverse itself, which can be non-trivial.
\end{complexitybox}

\begin{takeawaybox}
Command packages an action plus its receiver and arguments into an object
with \texttt{execute}/\texttt{undo}, unlocking undo stacks, queues, and
macros. The interview keyword is almost always ``undo/redo.'' It is the
backbone of editors and task schedulers.
\end{takeawaybox}
